\subsection*{II.}
在证明约化定理之前，我们先说明形式的线性族的一些性质。这里所谓\emph{$K$ 上形式的线性族}，是指一个由次数相同的形式组成、并对下述运算封闭的系统：只要系统含有 $\theta_1$ 和 $\theta_2$，它就总也含有 $c_1\theta_1+c_2\theta_2$；其中 $c_1,c_2$ 是给定有理性域 $K$ 中的量，而 $K$ 包含各 $\theta$ 的系数域。在这些形式中，总可以找到有限多个——设为 $\rho$ 个——在 $K$ 上线性无关的形式，而任意 $\rho+1$ 个形式都满足一个系数取自 $K$ 的线性关系；这个族的\emph{秩为 $\rho$}（相对于 $K$）。\NoetherSrcNote{*)}{若去掉次数相同这一限制，或者去掉 $K$ 包含各 $\theta$ 的系数域这一限制，就可得到更一般的形式的线性族；它们的秩不必有限。}有如下容易看出的\emph{事实}：若 $\mathcal L$ 和 $\mathcal T$ 是 $K$ 上秩同为 $\rho$ 的两个线性族，而 $\mathcal T$ 是 $\mathcal L$ 的子族，则 $\mathcal L$ 与 $\mathcal T$ 相同（等价）。

有了这些准备，我们转向约化定理。I 中的恒等式
\[
  \theta=\sum PZ
\]
对上述恒等式两边施以同一极化过程，就得到关于 $P\theta$ 的相应恒等式；这个过程记为 $P$。所以约化定理也就同时给出了\emph{$\theta$ 的一切极化式}由特殊极化式 $PZ$ 之和表示的方法。另一方面，这些特殊极化式显然包含在全部极化式之中，\emph{因此约化定理所断言的，正是这两个形式的线性族等价}；特别地，它也断言如下两个子族等价：其中每个形式在每一个\emph{单独的}变量组中都与 $\theta$ 具有相同次数。为了证明这种等价性，我们再插入一个由形式 $Z$ 所决定的中间族。为简单起见，\emph{先证明三个二元变量组 $(N=3,n=2)$ 的情形}，然后再简要说明与之平行的一般证明。

于是，设 $\theta(x,y,z)$ 在三个变量组中的次数分别为 $\alpha,\beta,p$；这三个变量组为
\[
  x_1x_2;\quad y_1y_2;\quad z_1z_2;
\]
暂且把 $\theta$ 的系数视为不定元 $a$（也可以视为有理数）。现在考虑以下三个形式的线性族：

1. 族 $\mathcal L$，它由所有形式 $f(x,y,z)$ 组成；这些形式由 $\theta$ 通过极化过程 $P$ 得到，且在各变量组 $x,y,z$ 中均与 $\theta$ 同次；特别地，$\mathcal L$ 也含有 $\theta$。把 $f(x,y,z)$ 看成关于 $x,y,z$ 及不定元 $a$ 的形式时，其系数域就是有理数域 $R$；$K$ 也必须取作 $R$，因为只有在 $\theta_1,\theta_2$ 取有理数值时，$c_1P_1+c_2P_2$ 才仍是一个过程 $P$；设 $\mathcal L$ 的秩为 $\rho$（相对于 $R$）。

2. 族 $\mathcal S$，它由所有按下式定义的形式 $g(\lambda;x,y)$ 组成：
\[
  g(\lambda;x,y)=f(x,y,\lambda_1x+\lambda_2y),
\]
或者用极化过程表示为：
\[
\begin{aligned}
  g(\lambda;x,y)
  &=\frac1{p!}\left[(\lambda_1x+\lambda_2y)\frac{\partial}{\partial z}\right]^p f(x,y,z)\srcfnmark{*)} \\
  &=g_0(xy)\lambda_1^p+g_1(x,y)\lambda_1^{p-1}\lambda_2+\cdots+g_p(xy)\lambda_2^p,
\end{aligned}
\]
其中有：
\[
  i!(p-i)!\,g_i(xy)=
  \left(y\frac{\partial}{\partial z}\right)^i
  \left(x\frac{\partial}{\partial z}\right)^{p-i}f(xyz).
\]
\srcfntext{*)}{其含义与 I 中相同：
\[
 \left(t\frac{\partial}{\partial x}\right)=t_1\frac{\partial}{\partial x_1}+t_2\frac{\partial}{\partial x_2},
\]
因而特别有：
\[
 \left[(\lambda_1x+\lambda_2y)\frac{\partial}{\partial z}\right]
 = (\lambda_1x_1+\lambda_2y_1)\frac{\partial}{\partial z_1}
 + (\lambda_1x_2+\lambda_2y_2)\frac{\partial}{\partial z_2},
\]
\[
 \left[z\left(\frac{\partial}{\partial\lambda_1}\frac{\partial}{\partial x}
 +\frac{\partial}{\partial\lambda_2}\frac{\partial}{\partial y}\right)\right]
 =z_1\left(\frac{\partial}{\partial\lambda_1}\frac{\partial}{\partial x_1}
 +\frac{\partial}{\partial\lambda_2}\frac{\partial}{\partial y_1}\right)
 +z_2\left(\frac{\partial}{\partial\lambda_1}\frac{\partial}{\partial x_2}
 +\frac{\partial}{\partial\lambda_2}\frac{\partial}{\partial y_2}\right).
\]}
这样，各 $g_i(xy)$ 就给出了恒等式 (1) 中的形式 $Z$。各 $g$ 是关于 $x,y,\lambda$ 及不定元 $a$、且系数为有理数的形式；设 $\mathcal S$ 的秩为 $\sigma$（相对于 $R$）。

3. 族 $\mathcal T$，它由 $\mathcal S$ 经过逆极化过程得到，正如 $\mathcal S$ 由 $\mathcal L$ 得到一样；形式 $h$ 属于 $\mathcal T$，并定义为：
\[
\begin{aligned}
 h(x,y,z)
 &=\frac1{p!}\left[z\left(\frac{\partial}{\partial\lambda_1}\frac{\partial}{\partial x}
 +\frac{\partial}{\partial\lambda_2}\frac{\partial}{\partial y}\right)\right]^p g(\lambda;xy)\\
 &=h_0(xyz)+h_1(xyz)+\cdots+h_p(xyz),
\end{aligned}
\]
其中由 2. 可知：
\[
  h_i(xyz)=
  \left(z\frac{\partial}{\partial x}\right)^{p-i}
  \left(z\frac{\partial}{\partial y}\right)^i g_i(xy).
\]
所以，各个 $h_i$，从而 $h$ 本身，都是形式 $PZ$ 或这类形式的和；设 $\mathcal T$ 的秩为 $\tau$（相对于 $R$）。

由各形式 $h(x,y,z)$（属于 $\mathcal T$）的构造可见，它们通过极化过程 $P$ 从 $\theta$ 得到，并且分别在变量组 $x,y,z$ 中与 $\theta$ 具有相同次数；\emph{因此，族 $\mathcal T$ 是 $\mathcal L$ 的一个子族}。根据上面所述的\emph{事实}，为了证明等价性，现在只需证明秩相等。在证明秩相等时，不再用到 $f(xyz)$ 的特殊结构——即它们都是同一个形式 $\theta$ 的极化式。

\emph{a) 为了比较 $\mathcal L$ 与 $\mathcal S$ 的秩，使用引理 a)：由 $f(x,y,\lambda_1x+\lambda_2y)=0$ 必然推出 $f(x,y,z)=0$。}这一点直接由三个二元变量组之间的恒等关系
\[
  (xy)z+(yz)x+(zx)y=0,
  \qquad \text{其中 }(xy)=(x_1y_2-x_2y_1),\ldots,
\]
得到；该关系蕴含：
\[
  f[x,y,(zy)x+(xz)y]=(xy)^p f(x,y,z).
\]
由 $f(x,y,\lambda_1x+\lambda_2y)=0$（关于 $\lambda$ 恒等为零），作代入 $\lambda_1=(zy)$、$\lambda_2=(xz)$，并注意 $(xy)\ne0$，便必然得到 $f(x,y,z)=0$。

根据这个引理，\emph{$\mathcal S$ 中的形式与 $\mathcal L$ 中的形式之间存在一一对应}。事实上，由
\[
  f_1(x,y,\lambda_1x+\lambda_2y)=g(\lambda;xy),\qquad
  f_2(x,y,\lambda_1x+\lambda_2y)=g(\lambda;xy)
\]
必然得到：
\[
  f_1(xyz)=f_2(xyz).
\]
而且，\emph{$\mathcal S$ 中的每个关系，在 $\mathcal L$ 中唯一对应的那些形式之间仍然成立（反之亦然）}。因为由
\[
  c_1g^{(1)}(\lambda;xy)+c_2g^{(2)}(\lambda;xy)+\cdots+c_rg^{(r)}(\lambda;xy)=0
\]
根据上述引理必然得到：
\[
  c_1f^{(1)}(x,y,z)+c_2f^{(2)}(x,y,z)+\cdots+c_rf^{(r)}(x,y,z)=0.
\]
\emph{由此便证明了 $\mathcal L$ 与 $\mathcal S$ 的秩相等；$\rho=\sigma$。}

\emph{b) 为了比较 $\mathcal S$ 与 $\mathcal T$ 的秩，相应地使用引理 b)：由 $h(x,y,z)=0$ 必然推出 $g(\lambda;xy)=0$。}为证明这一点，注意根据 3.，$h(x,y,z)=0$ 等价于下列各幂乘积均为零：
\[
 \left(\frac{\partial}{\partial\lambda_1}\frac{\partial}{\partial x_1}
 +\frac{\partial}{\partial\lambda_2}\frac{\partial}{\partial y_1}\right)^{p_1}
 \left(\frac{\partial}{\partial\lambda_1}\frac{\partial}{\partial x_2}
 +\frac{\partial}{\partial\lambda_2}\frac{\partial}{\partial y_2}\right)^{p_2}g(\lambda;x,y),
 \qquad p_1+p_2=p,
\]
进一步微分可知，下列各幂乘积也都为零：
\[
\begin{aligned}
&\left(\frac{\partial}{\partial x_1}\right)^{\alpha_1}
 \left(\frac{\partial}{\partial x_2}\right)^{\alpha_2}
 \left(\frac{\partial}{\partial y_1}\right)^{\beta_1}
 \left(\frac{\partial}{\partial y_2}\right)^{\beta_2} \\
&\quad\cdot
 \left(\frac{\partial}{\partial\lambda_1}\frac{\partial}{\partial x_1}
 +\frac{\partial}{\partial\lambda_2}\frac{\partial}{\partial y_1}\right)^{p_1}
 \left(\frac{\partial}{\partial\lambda_1}\frac{\partial}{\partial x_2}
 +\frac{\partial}{\partial\lambda_2}\frac{\partial}{\partial y_2}\right)^{p_2}
 g(\lambda;xy)
\end{aligned}
\]
因此，这些幂乘积的任意线性组合也都为零，从而每个形式
\[
  f\left(\frac{\partial}{\partial x},\frac{\partial}{\partial y},
  \frac{\partial}{\partial\lambda_1}\frac{\partial}{\partial x}
  +\frac{\partial}{\partial\lambda_2}\frac{\partial}{\partial y}\right)g(\lambda,xy).
\]
也都为零。特别选取 $f$，使得
\[
  f(x,y,\lambda_1x+\lambda_2y)=g(\lambda;xy),
\]
于是也得到：
\[
  g\left(\frac{\partial}{\partial\lambda};
    \frac{\partial}{\partial x},
    \frac{\partial}{\partial y}\right)g(\lambda,x,y)=0;
\]
或者，若令
\[
  g(\lambda,xy)=\sum G_i\lambda_1^{\varkappa_1}\lambda_2^{\varkappa_2}
  x_1^{\varkappa_3}x_2^{\varkappa_4}y_1^{\varkappa_5}y_2^{\varkappa_6},
\]
其中各 $G_i$ 是关于不定元 $a$ 的有理系数线性形式，则有
\[
  \sum \varkappa_1!\varkappa_2!\cdots\varkappa_6!\,G_i^2=0
\]
从而 $g(\lambda;xy)=0$。

由引理 b)，完全仿照 a) 即得：\emph{$\mathcal S$ 与 $\mathcal T$ 的秩相等；$\sigma=\tau$。}

所以，由 a) 和 b) 可得 $\rho=\tau$；又因为 $\mathcal T$ 是 $\mathcal L$ 的子族，这两个线性族的等价性就得到了证明。于是，$\mathcal L$ 中的每个形式，特别是 $\theta$，都可以表示成 $\rho$ 个线性无关形式 $h(x,y,z)$ 的有理系数线性组合：
\[
  \theta=\sum c_i h^{(i)}(x,y,z)=\sum PZ.
\]
这个对不定系数 $a$（即 $\theta$ 的系数）所推得的恒等式，在用任意有理性域中的量代替不定元 $a$ 时仍然成立；\emph{因此，三个二元变量组情形下的约化定理已经得到一般证明。}

现在来看一般证明：对于 $N$ 个变量组、每组各含 $n$ 个变量的情形，证明完全平行。设 $\theta(A)$ 的次数分别为 $\alpha_i$，相应的变量组为 $A_i$。我们再次考虑以下三个形式的线性族：

1. 族 $\mathcal L$，由所有形式 $f(A)$ 组成；这些形式由 $\theta(A)$ 通过极化过程 $P$ 得到，并且在每一个单独的变量组 $A_i$ 中都与 $\theta(A)$ 同次；

2. 族 $\mathcal S$，由所有形式 $g(\lambda;A_1\cdots A_n)$ 组成；这些形式由 $f(A)$ 经过如下代入得到：
\[
\begin{aligned}
 A_{n+1}&=\lambda_{n+1}^{(1)}A_1+\lambda_{n+1}^{(2)}A_2+
        \cdots+\lambda_{n+1}^{(n)}A_n,\\
 &\vdotswithin{=}\\
 A_N&=\lambda_N^{(1)}A_1+\lambda_N^{(2)}A_2+
        \cdots+\lambda_N^{(n)}A_n
\end{aligned}
\]
这里，各个 $\lambda$ 幂乘积在 $g(\lambda;A_1\cdots A_n)$ 中的系数给出了恒等式 (1) 中的形式 $Z$；

3. 族 $\mathcal T$，由所有按下式定义的形式 $h(A)$ 组成：
\[
\begin{aligned}
 h(A)=&\left[A_{n+1}\left(
 \frac{\partial}{\partial\lambda_{n+1}^{(1)}}\frac{\partial}{\partial A_1}
 +\cdots+
 \frac{\partial}{\partial\lambda_{n+1}^{(n)}}\frac{\partial}{\partial A_n}\right)\right]^{\alpha_{n+1}} \\
&\cdots
\left[A_N\left(
 \frac{\partial}{\partial\lambda_N^{(1)}}\frac{\partial}{\partial A_1}
 +\cdots+
 \frac{\partial}{\partial\lambda_N^{(n)}}\frac{\partial}{\partial A_n}\right)\right]^{\alpha_N}
 g(\lambda;A_1\cdots A_n);
\end{aligned}
\]
$h(A)$ 是形式 $PZ$ 的和。

由于每 $n+1$ 个变量组之间都存在恒等关系，引理 a) 仍然成立；引理 b) 也同样成立，因为其证明根本没有用到变量的数目。因此，$\mathcal L$ 与 $\mathcal T$ 的秩相等；又因为 $\mathcal T$ 是 $\mathcal L$ 的子族，二者便相同。所以恒等式
\[
  \theta=\sum PZ
\]
对于不定系数成立，因而当系数取任意有理性域中的量时也成立；\emph{至此，约化定理的一般情形也得到了证明。}
